% Secao - mobilidade medica (lado da oferta). Numeros de 15/16/17.
% Disciplina: descritivo, associativo. % src aponta a script/CSV.

\section{O lado da oferta: o deserto de fluxo tem dois lados}
\label{sec:medicos}

A Seção~\ref{sec:divergencia} deixou um enigma. Os municípios perto-mas-isolados
não são desprovidos de hospital --- têm a maior presença de leito local entre os
quadrantes de deserto --- e ainda assim 90\% de suas internações escapam para
hubs distantes. O leito existe; o cuidado, não. A hipótese natural é que falte ao
hospital local o insumo que retém o paciente: o médico, e em particular o
especialista. Esta seção leva a hipótese ao cadastro de força de trabalho do SUS,
e vai além do estoque: pergunta se os médicos também \emph{se movem}, ao longo do
tempo, na direção em que os pacientes já viajam.

Usamos o cadastro de profissionais do CNES (CNES-PF), vínculo médico $\times$
estabelecimento $\times$ mês, 2015--2023 --- 128,9 milhões de vínculos médicos
---, identificando o médico pela família CBO de medicina (verificada contra o
registro no CRM) e deduplicando-o pelo Cartão Nacional de Saúde, presente em
100\% dos registros.% src: 15/16; CNS_PROF; CBO 2231/2251/2252/2253 = CONSELHO 71
A cada médico, em cada ano, atribuímos o \emph{município primário} --- onde
concentra mais horas ---, e registramos como migração toda mudança de município
primário entre anos consecutivos.

\subsection{Estoque: menos médicos, menos especialistas, apesar do leito}

A Tabela~\ref{tab:medicos} é direta. Os municípios perto-mas-isolados têm 1,48
médico por mil habitantes contra 2,22 no quadrante conectado --- uma oferta uma
vez e meia menor --- e o déficit é mais agudo onde mais importa: especialistas
(famílias cirúrgica e diagnóstico-terapêutica) são 12,3\% do corpo clínico nos
desertos de fluxo contra 18,6\% no quadrante conectado, e apenas 5,4\% no deserto
pleno.% src: physician_supply_panel / bloco especialista por quadrante
É a peça que faltava ao enigma do leito: o hospital local existe, mas sem a
densidade médica --- e sobretudo sem o especialista --- que faz o paciente
permanecer. A decomposição por tipo de estabelecimento afia o ponto. A atenção
primária é semelhante em todos os quadrantes (0,80 médico de unidade básica por
mil habitantes nos perto-mas-isolados contra 0,92 no conectado), mas o médico
hospitalar é escasso (0,54 contra 0,72) e o especialista ainda mais (0,25 contra
0,51, e apenas 0,11 no deserto pleno).% src: 19_physician_setting / physician_setting_by_quadrant.csv
O deserto de fluxo não carece do piso da atenção básica --- universalizado pela
estratégia de saúde da família ---, mas do andar hospitalar e especializado,
exatamente o que retém o caso complexo. A oferta não substituída pelo leito é
a que o fluxo revela escapando.

\begin{table}[t]
\centering
\caption{Oferta e mobilidade médica por quadrante de deserto, 2015--2023}
\label{tab:medicos}
\footnotesize
\begin{tabular}{lrrrr}
\toprule
Quadrante & Médicos/ & \% espe- & Retenção & Saldo migr. \\
          & 1.000    & cialista & anual    & /1.000 méd. \\
\midrule
Perto, deserto em fluxo & 1{,}48 & 12{,}3 & 0{,}77 & $-184$ \\
Deserto pleno           & 1{,}13 & 5{,}4  & 0{,}77 & $-351$ \\
Longe, conectado        & 1{,}41 & 10{,}5 & 0{,}78 & $-250$ \\
Sem deserto             & 2{,}22 & 18{,}6 & 0{,}81 & $+2$ \\
\bottomrule
\end{tabular}
\par\smallskip
\begin{minipage}{0.82\linewidth}
\footnotesize
\textit{Nota:} Medianas municipais, 2015--2023. Especialista $=$ famílias CBO
cirúrgica e diagnóstico-terapêutica. Retenção $=$ fração dos médicos primários em
$t$ ainda primários no município em $t+1$. Saldo migratório $=$ (entradas $-$
saídas) por mil médicos. Fonte: CNES-PF, IBGE.
\end{minipage}
\end{table}

\subsection{Dinâmica: os desertos sangram médicos}

A oferta menor não é um estoque estático --- é alimentada por um fluxo de saída.
A retenção anual de médicos é menor nos desertos (0,77 contra 0,81 no quadrante
conectado), e o saldo migratório separa nitidamente os dois mundos: os três
quadrantes de deserto perdem médicos ano após ano --- $-184$ por mil médicos nos
perto-mas-isolados, $-351$ no deserto pleno --- enquanto o quadrante conectado é
o único que \emph{ganha} ($+2$).% src: physician_by_quadrant.csv; net_por_1000_med
O deserto de fluxo, portanto, tem dois lados que apontam na mesma direção: o
paciente viaja para fora, e o médico também migra para fora.

\subsection{As duas fugas correm juntas --- mas não para o mesmo hub}

A pergunta seguinte é se os dois drenos têm o mesmo destino. A resposta é
parcial, e a nuance importa. As duas redes --- migração de médicos e fluxo de
pacientes --- são moderadamente correlacionadas no nível das arestas
município$\rightarrow$município (correlação de 0,44 sobre 54 mil pares comuns;
Figura~\ref{fig:redes_medicas}): médicos e pacientes deixam os mesmos lugares por
corredores parecidos.% src: physician_vs_patient_network.csv; r_log 0.455
Mas, \emph{dentro} dos desertos de fluxo, os destinos divergem. O principal
destino dos médicos que saem coincide com o principal hub de pacientes em apenas
41\% dos municípios perto-mas-isolados, contra 53\% nos conectados, e somente
16\% dos médicos que partem rumam para o hub que recebe a maior parte dos
pacientes do município.% src: tophub_match_by_quadrant.csv
O padrão é coerente: o paciente busca o centro de referência mais próximo que
resolva seu caso; o médico que deixa o deserto tende a ir mais longe, com
frequência à capital, por razões de carreira que não são as do paciente. As duas
fugas esvaziam o deserto, mas não o reabastecem no mesmo ponto.

\subsection{O que esta seção contribui (e o que não)}

O deserto de fluxo não é só um fenômeno de demanda. Ele é duplo: os municípios
que perdem pacientes para hubs distantes são também os que têm menos médicos,
menos especialistas e saldo migratório médico negativo --- um arranjo
auto-coerente em que a ausência do cuidado e a ausência de quem o presta
caminham juntas. A leitura permanece associativa. A concentração de especialistas
em centros de referência é, em parte, organização racional do sistema, e nada
aqui identifica se a saída dos médicos \emph{causa} o esvaziamento do hospital
local ou apenas o acompanha. Duas limitações de medida acompanham o resultado: o
CNES capta vínculo de trabalho, não residência --- o campo de município de
residência do profissional está preenchido em apenas 7\% dos vínculos, o que
impede separar deslocamento pendular de migração efetiva ---;% src: UFMUNRES fill-rate 7.3%
e a oferta é medida por presença de vínculo, não por horas efetivamente
ofertadas no SUS. O custo em saúde desse arranjo de oferta, como o da demanda na
Seção~\ref{sec:mortalidade}, fica reservado ao desenho causal.
